\section{Renormalization in auxiliary field approach}\label{sec:renorm}

The renormalization of the effective Lagrangian Eq.~(\ref{efflag}) is in complete analogy to that with a fundamental auxiliary field~\cite{Bagan:1993zv,Ji:2017oey}, and suffices to render Green's functions of the elementary fields finite. However, extra complications arise when renormalizing the LCOs $J_1, {\overline J}_1$ in Eq.~(\ref{Eq:repl}).
We first consider the perturbative renormalization in dimensional regularization (DR) in a covariant gauge. In such a case, gauge-invariant LCOs can, in general, mix with operators of the same or lower mass dimension under renormalization. The mixing operators can be of the following three types: 1) gauge-invariant operators, 2) BRST exact operators,
3) operators that vanish by equation of motion (EOM)~\cite{Collins:1984xc}. For the LCO formed by a quark and a fundamental auxiliary field, such a mixing does not occur because the operator appearing there is already of the lowest dimension and, if chiral symmetry is preserved, the mixing between operators of the same dimension is forbidden~\cite{Chen:2017mie}. This is no longer true in the present case.

In DR, the operators allowed by BRST symmetry to mix with $J_1^{\mu\nu}$ are
\begin{align}\label{mixingop}
%J_1^{\mu\nu}&=F^{\mu\nu}_a\mathcal Q_a, \non\\
J_2^{\mu\nu}&=n_{\rho}(F^{\mu\rho}_a n^\nu-F^{\nu\rho}_a n^{\mu}){\mathcal Q}_a/n^2,\non\\
%J_3^{\mu\nu}&=(-in^\mu A_a^{\nu}+in^\nu A_a^\mu)(in\cdot D \mathcal Q)_a/n^2,
J_3^{\mu\nu}&=(-in^\mu A_a^{\nu}+in^\nu A_a^\mu)((in\cdot D-m) \mathcal Q)_a/n^2,
\end{align}
where $J_2^{\mu\nu}$ is gauge invariant, $J_3^{\mu\nu}$ is proportional to the EOM of $\mathcal Q$ and therefore vanishes in a physical matrix element (for $m=0$ the above mixing pattern has been given in~\cite{Dorn:1981wa}).

In the presence of mixing, the renormalized operators can be written in a triangular mixing form
\begin{align}\label{oprenormmix}
\begin{pmatrix}
 J_{1, R}^{\mu\nu} \\  J_{2, R}^{\mu\nu} \\ J_{3, R}^{\mu\nu}
\end{pmatrix}
&=
\begin{pmatrix}
Z_{11} & Z_{12} & Z_{13}\\ 0 & Z_{22} & Z_{23} \\ 0 & 0 & Z_{33}
\end{pmatrix}
\begin{pmatrix}
 J_1^{\mu\nu}\\  J_2^{\mu\nu} \\ J_3^{\mu\nu}
\end{pmatrix}.
\end{align}
The renormalization constants in Eq.~(\ref{oprenormmix}) are not all independent. When $\nu=z$, $J_2^{z\mu}$ becomes degenerate with $J_1^{z\mu}$,
therefore they must have the same renormalization, implying
\beq\label{Zrelation}
Z_{11}+Z_{12}=Z_{22}, \hspace{5em} Z_{13}=Z_{23}.
\eeq
This is consistent with explicit one-loop calculations~\cite{Dorn:1981wa}.
The above result indicates that the individual components of $J_1^{\mu\nu}$ renormalize independently with the following simplified renormalization equations:
\begin{align}\label{oprenormmixsim2x2}
\begin{pmatrix}
 J_{1, R}^{z\mu} \\ J_{3, R}^{z\mu}
\end{pmatrix}
&=
\begin{pmatrix}
Z_{22} & Z_{13} \\  0 & Z_{33}
\end{pmatrix}
\begin{pmatrix}
 J_1^{z\mu} \\ J_3^{z\mu}
\end{pmatrix};\hspace{1em}
J_{1, R}^{ti}=Z_{11} J_1^{ti}.
\end{align}
$J_{1}^{ij}$ shares the same renormalization with $J_{1}^{ti}$. The reason that
($J_{1}^{ti}$, $J_{1}^{ij}$) and $J_1^{z\mu}$ operators have different renormalizations is because the presence of a four-vector $n^\mu$ breaks Lorentz symmetry. Note that $J_2^{\mu i}$ is not independent, and henceforth will be neglected.

Our calculation show that no linear divergences occur in the one-loop correction to $J_1^{\mu\nu}$. Actually, the power divergence structure of this operator has been considered in Refs.~\cite{Wang:2017qyg,Wang:2017eel}, but using a simple cutoff regularization. One has to be cautious when dealing with power divergences in cutoff regularization, since it in general breaks gauge invariance in QCD (except for the lattice cutoff), and might obscure the structure of genuine power divergences of the theory. To avoid this, we have chosen to work in DR and kept track of linear divergences by expanding around $d=3$, as they appear as poles at $d=3$. In this way, we can extract linear divergences gauge invariantly and confirm that no linear divergences occur in the one-loop correction to $J_1^{\mu\nu}$. The same conclusion can actually be reached in Ref.~\cite{Wang:2017eel} too, if a gauge-invariant regulator is used. We have explicitly checked this following the procedure outlined above.

In Eq.~(\ref{efflag}), we also introduced a ``residual mass" term for the auxiliary field $\mathcal Q$. The reason is to keep track of potential mass renormalization, which might be absent in DR when expanded at $d=4$, but nevertheless appears in a cutoff regularization such as lattice regularization. In a cutoff regularization, an effective mass for the auxiliary ``heavy quark" will be generated by radiative corrections even if it does not exist at leading order.
This is indeed what happens in our present case, and we have $m=\delta m\sim{\cal O}(1/a)$ with $a$ being the inverse cutoff or lattice spacing~\cite{Maiani:1991az}. In perturbation theory, $\delta m$ appears from $O(\alpha_s)$.
Such a mass term serves the purpose of absorbing power divergences arising from the Wilson line self energy when integrating out the auxiliary field, as will be seen below. Apart from this, there is no other power divergence in the theory. Moreover, for dimensional reasons, there is no other antisymmetric operator that can mix with $J_i^{\mu\nu}$ discussed previously. Therefore in a gauge-invariant cutoff scheme, the operator renormalization remains the same as in DR.

In Ref.~\cite{Ji:2017oey} it was shown that BRST invariance of the Lagrangian requires a dependence of $m$ on the signature of $n$~\cite{Dorn:1986dt}, which yields a vanishing mass for a lightlike $n^\mu$ and a real mass for a timelike $n^\mu$. When $n^\mu$ is spacelike, $m=\delta m=i\overline{\delta { m}}$ is imaginary. The value of $\overline{\delta { m}}$ at $O(\alpha_s)$ can be easily obtained from Ref.~\cite{Ji:2017oey} as $\overline{\delta { m}}={\pi\alpha_s C_A}/(2a)$.
